\SourcePage{348}{353}
\section{Crossing invariance}

Representing the scattering amplitudes $M_{fi}$ by Feynman integrals reveals a remarkable symmetry, as follows.

Without changing its direction, any incoming external line of a Feynman diagram may represent either an initial particle or a final antiparticle; any outgoing line may represent either a final particle or an initial antiparticle. Passing from particle to antiparticle simultaneously changes the meaning of the 4-momentum $p$ assigned to the line: $p=p_{\text{e}}$ for a particle, say an electron, and $p=-p_{\text{p}}$ for a positron. The assigned polarization also changes. An incoming line carries the wave amplitude $u$ and an outgoing line $u^*$; $u=u_{\text{e}}$ for an electron, but $u=u_{\text{p}}^*$ for a positron. Replacing $u$ by $u^*$ reverses the particle's spin projection, or helicity.

For a photon, a genuinely neutral particle, changing the meaning of an external line merely interchanges emission and absorption. A photon line of momentum $k$ represents either absorption with $k_{\text{abs}}=k$, or emission with $k_{\text{em}}=-k$ and the opposite helicity.

Changing the meaning of the external lines in this way amounts to passing from one crossed reaction channel to the others. Hence the same amplitude, as a function of the momenta at the diagrams' free ends, describes every reac\SourcePageJoin{349}{354}tion channel\footnote{If a channel is forbidden by 4-momentum conservation, the delta function occurring as a common factor in (64.5) automatically makes its transition probability zero.}. Only the interpretation of its arguments changes: replacing a particle by an antiparticle means substituting $p_i\to-p_f$, where $p_i$ is the initial-particle 4-momentum in one channel and $p_f$ the final-particle 4-momentum in another. This property of a scattering amplitude is called \emph{crossing symmetry}, or \emph{crossing invariance}.

For the invariant amplitudes introduced in \S~70 as functions of kinematic invariants, the functions are the same in every channel, but their arguments range over a different physical region in each. Thus Feynman integrals define invariant amplitudes as analytic functions; values in different physical regions are analytic continuations of the function specified in one region. Since their integrands contain singularities, the invariant amplitudes have singularities determined by the integral expressions, including the pole-bypass prescription. Once calculated from Feynman integrals in one channel, their analytic continuation to other channels automatically incorporates these singularities.

Crossing invariance is more than the scattering-matrix properties implied by general space-time symmetry. Those requirements equate amplitudes for processes related by interchanging initial and final states and replacing all particles by antiparticles, with every particle momentum $\mathbf p$ unchanged and the signs of the angular-momentum projections reversed. This is $CPT$ invariance\footnote{Formally passing between these reactions by reversing all 4-momenta in the Feynman diagrams corresponds to $CPT$ as four-dimensional inversion.}. Crossing invariance permits this transformation not only for all particles together, but for any one particle separately.
